\section{Parameter choices and the final contradiction}
\label{sec:parameters}

\begin{lemma}[The comparison gap]\label{lem:gap}
Let \(\kappa\ge1\) and \(\nu>2\kappa\). There are positive
rational numbers \(\theta,A,B,C\) and a real \(\eta\in(0,1)\)
such that
\begin{gather}
 0<\theta<A<B<1,\qquad C>1,\qquad
 C\theta<1,\quad CB<1,\quad C\theta/B>1,\quad B/A>1,
 \label{eq:parameter-ratios}\\
 g:=\nu\bigl(A(1-\eta)-\theta\bigr)-\kappa(1-\theta)>0.
 \label{eq:gap}
\end{gather}
\end{lemma}
\begin{proof}
Choose a rational \(b_0\) with
\[
 \frac12<b_0<1-\frac{\kappa}{\nu}.
\]
For a sufficiently small positive rational \(\delta\), put
\(\theta=1-\delta\), \(A=1-b_0\delta\). Then \(0<\theta<A<1\),
\[
 \nu(A-\theta)-\kappa(1-\theta)
     =\bigl(\nu(1-b_0)-\kappa\bigr)\delta>0,
\]
and
\(\theta-A^2=(2b_0-1)\delta-b_0^2\delta^2>0\).
Choose rationals
\[
 A/\theta<C<1/A,\qquad A<B<\min(1/C,C\theta).
\]
The intervals are nonempty because \(A^2<\theta\).
Also \(C\theta<\theta/A<1\), giving all the ratio inequalities.
Finally choose \(\eta>0\) small enough to preserve the positive gap.
\end{proof}

The constants in the preceding estimates can all be made small in the
required order. To see this without suppressing any number-field cost,
put \(E=E_{\rm ar}+E_{\rm an}\) and define
\[
 \begin{split}
 \mathcal C&=R_0+\kappa\log D_\xi+
                    \kappa(\log\Xi+\log(3/2)),\\
 \mathcal W&=\kappa\theta+\log4+\log(2K\Gamma)+\nu
       +\log(2R_0K)+\kappa\log(2(KB_0+2)).
 \end{split}
\]
All these quantities are nonnegative. The exact error identity is
\begin{equation}\label{eq:error-ledger}
 \begin{split}
 E={}&\frac{\nu}{F_0}
   +\frac{\kappa\Lambda F_0m+(\kappa+2)\log2}{v_0}
   +\mathcal C\frac K{w_0}\\
   &+\kappa(\Lambda+\log D_\gamma)\sum_{i=1}^m\frac1{w_i}
     +\frac{\mathcal W}{w_*}.
 \end{split}
\end{equation}
In particular, algebraic denominators and unconstrained conjugates
appear explicitly rather than being hidden in the \(o(1)\) in \(H\).

\begin{proposition}[Admissible finite choices]\label{prop:choices}
Assuming \eqref{eq:bad-approximations}, all fixed data can be selected
so that the interpolation theorem applies and
\begin{equation}\label{eq:margins}
 E<g,\qquad cL>\kappa+E.
\end{equation}
\end{proposition}
\begin{proof}
Choose \(\theta,A,B,C,\eta\) by Lemma~\ref{lem:gap}, and set
\(\epsilon_0=\min(g/2,1/2)>0\).
First choose \(F_0>2/\theta\) such that
\(\nu/F_0<\epsilon_0/3\).
For an integer \(m\ge1\), define
\begin{equation}\label{eq:dimension-parameters}
 K=\lfloor C^m\rfloor,\qquad w_0=B^{-m},\qquad
 v_0=2K\theta^m w_0.
\end{equation}
They satisfy
\[
 K(w_0/v_0)\theta^m=\tfrac12,\qquad
 K\theta^m\le(C\theta)^m<1.
\]
The second inequality allows the common-fiber branch as well as the
distinct-center branch.
Moreover
\[
 \frac K{w_0}\le(CB)^m\longrightarrow0,\qquad
 v_0\sim2(C\theta/B)^m\longrightarrow\infty,\qquad
 \frac m{v_0}\longrightarrow0.
\]
Substitution in \eqref{eq:collision-limit} gives the exact expression
\begin{equation}\label{eq:L-dimension}
 L=\frac{\eta^2}{2(m+1)}(B/A)^m\longrightarrow\infty.
\end{equation}
Thus one can fix a finite \(m\) large enough that
\[
 \frac{\kappa\Lambda F_0m+(\kappa+2)\log2}{v_0}
       +\mathcal C\frac K{w_0}<\epsilon_0/3,
 \qquad cL>\kappa+1.
\]
This fixes \(m,K,w_0,v_0\), and hence \(\mathcal W\).

Now choose \(\sigma_0\) and the weight thresholds for
Theorem~\ref{thm:interpolation}, imposing also its common-fiber
volume inequality. Choose the finitely many approximants
\(p_i/q_i\) successively from \eqref{eq:bad-approximations}.
Their weights can satisfy every separation threshold and can all
exceed a common bound large enough that
\[
 \kappa(\Lambda+\log D_\gamma)\sum_i\frac1{w_i}
               +\frac{\mathcal W}{w_*}<\epsilon_0/3.
\]
Increase the denominator thresholds as needed to ensure \(q_i\ge2\)
and \(p_i\ne0\).
This step is valid because the interpolation weight thresholds are
independent of the centers, which themselves involve \(p_i/q_i\).

Equation~\eqref{eq:error-ledger} now gives \(E<\epsilon_0<g\);
also \(cL>\kappa+1>\kappa+E\).
Every choice is finite. Only after fixing this entire list do we let
the sufficiently divisible parameter \(H\) increase.
\end{proof}

\begin{proof}[Proof of Theorem~\ref{thm:compatible}]
Normalize \((\gamma,\xi)\) by Lemma~\ref{lem:torsion}.
Fix \(\nu>2\kappa\) and suppose \eqref{eq:bad-approximations}.
Make the choices of Proposition~\ref{prop:choices}.
For arbitrarily large divisible \(H\), interpolation supplies a
nonzero selected determinant. Propositions~\ref{prop:arithmetic} and
\ref{prop:analytic} then give
\[
 -\kappa(1-\bar b_H)-E_{\rm ar}
 \le E_{\rm an}+o(1)
       +\max\{-cL_H,-\nu(A(1-\eta)-\bar b_H)\}.
\]
Both entries of this maximum are incompatible with the lower bound
for sufficiently large \(H\).
For the collision entry, \(0\le\bar b_H\) implies
\(\kappa(1-\bar b_H)\le\kappa\), whereas
\(cL_H\to cL>\kappa+E\).
For the other entry, since \(\nu>\kappa\) and
\(\bar b_H\le\theta\),
\[
 \nu\bigl(A(1-\eta)-\bar b_H\bigr)
     -\kappa(1-\bar b_H)
 \ge \nu\bigl(A(1-\eta)-\theta\bigr)-\kappa(1-\theta)
 =g>E.
\]
The strict margins absorb the common remainder \(o(1)\).
This contradiction excludes \eqref{eq:bad-approximations}.
It proves \(\ELB(x,2\kappa)\), including the bound for unreduced
rational representations. Lemma~\ref{lem:endpoint} proves
irrationality and the exponent bounds.
The rational-logarithm theorem and algebraic-base corollary follow
from the specializations already given.
\end{proof}

\begin{remark}[What the exponent threshold means]
The condition \(\nu>2\kappa\) is used to choose
\(1/2<b_0<1-\kappa/\nu\).
It supplies the two simultaneous requirements \(A^2<\theta\) and a
positive determinant gap. This is the precise source of \(2d/s\).
At the boundary \(\nu=2\kappa\), this argument gives neither that
interval nor a strict gap.
For rational logarithms, exponent two means the eventual inequalities
for every \(2+\varepsilon\); it does not assert an eventual lower
bound with the exact exponent \(2\).
\end{remark}
